\begin{tcolorbox}[colback=white,colframe=genderMasculine,title={Table colours},fonttitle=\bfseries]
The background highlights common object-case patterns in the meanings shown.
An object can be omitted in context; a colour does not mean that every sentence
must contain one.

\smallskip
\begin{tabularx}{\linewidth}{@{}lX@{}}
\colorbox{caseAccusativeBg}{\strut Blue: Akk} &
Accusative object: \textit{Ich besuche den Freund.}\\
\colorbox{caseDativeBg}{\strut Green: Dat} &
Dative object: \textit{Ich helfe dem Freund.}\\
\colorbox{caseAccusativeDativeBg}{\strut Purple: Akk and Dat} &
Both cases, often together: \textit{Ich gebe dem Freund das Buch.}
They can also occur in different meanings, as with \textit{raten} and \textit{schmecken}.\\
\colorbox{black!5!white}{\strut Grey} &
No Akk/Dat object pattern highlighted in the usual meaning shown.\\
\end{tabularx}

\smallskip
For reflexive entries, the pronoun's case is included:
\textit{ich erinnere mich} is blue;
\textit{ich merke mir etwas} is purple.
\textbf{KONSTRUKTIONEN} lists object-case and prepositional patterns.
Cases governed by a preposition are separate from the background classification:
\textit{auf den Brief reagieren} does not make \textit{reagieren} blue.
An extra beneficiary (\textit{Ich kaufe dir ein Buch}) does not by itself
change a blue entry to purple. The colours are a learning guide, not a list of
every possible construction or meaning. Pink marks the nominative complement
of \textit{sein/werden}; yellow identifies modal verbs.
Compact separable compounds share their stem's box, not necessarily its case pattern.
\end{tcolorbox}

\begin{tcolorbox}[colback=referenceYellowBg,colframe=genderPlural]
\textbf{Table convention.} A single em dash (---) under \textbf{IMPERATIV}
means that no normal imperative is shown. A dash in a person row means that
this person is not normally used in the stated meaning: for example,
\textit{geschehen} (to happen) is shown with \textit{es/sie}, not \textit{ich/du}.
Under \textbf{KONSTRUKTIONEN} or
\textbf{TRENNBARE VERBEN}, it means that the table intentionally lists no item
in that category; it is not an accidentally blank section. Separable compounds
are listed compactly beneath their base verb, with a vertical bar marking the
split point (for example, \textbf{aus$|$stellen}). Each compound remains in the
German and English indexes, which lead to the applicable stem table.
Stem forms help with conjugation, but check each compound's meaning, reflexive
pronoun and applicable persons. A compound can also select a different
auxiliary; for example,
\textbf{ab$|$biegen} forms the Perfekt with \textit{sein}:
\textit{ist abgebogen}.
\end{tcolorbox}
